\section{Background}

\subsection{Agentic Systems and Inter-Agent Communication}

% Sources: README.md, Abstract, Agent Communication, and Experiments;
% paper/PAPER_STORY.md, Section 1; paper/sections/01_introduction.tex.
Large language model (LLM)-based agents can interact with other agents to
coordinate work and delegate tasks on behalf of users. Such interaction allows
a task to involve multiple participants, each contributing information or
performing part of the work. Communication may proceed autonomously, with
limited human involvement in individual exchanges. In this setting,
delegation describes assigning work to another agent; the permission to
contact that agent is governed by the communication system's authorization
policy.

Inter-agent communication therefore involves both task messages and mechanisms
for governing their exchange. Agent identities identify the participants,
secure transport protects communication, and authorization policies specify
which contacts are permitted. We call the participant requesting contact the
\emph{initiating agent} and the participant receiving that request the
\emph{receiving agent}.

SAGA provides a concrete setting for these roles. Users register their agents
with a Provider that maintains agent contact information and user-defined
access-control policies. The Provider assists agents in enforcing these
policies, while access-control tokens mediate subsequent communication.
This separates the user's configuration of contact permissions from the
agents' autonomous exchange of task messages.

\subsection{Authentication, Authorization, and Confidentiality}

% Sources: paper/PAPER_STORY.md, Section 1;
% paper/sections/03_problem_formulation.tex, Authentication Is Not Authorization.
\emph{Authentication} answers ``Who is communicating?'' It establishes the
identity of a participant in a protocol exchange. Identity credentials and
checks on protocol messages provide the basis for associating an exchange
with that participant.

\emph{Authorization} answers ``Who is allowed to communicate?'' It determines
whether a participant has permission to perform the requested interaction
under an applicable policy. For contact authorization, the relevant decision
concerns an initiating agent's permission to communicate with a receiving
agent according to the receiving user's rules.

\emph{Confidentiality} answers ``Who can learn the exchanged information?''
It concerns restricting disclosure of protected information to the parties
permitted to learn it within the specified security model. Protected
information may include communication content and credentials exchanged by
the protocol.

These properties are complementary. Authentication supplies the participant
identity to which a policy can apply; authorization determines the permitted
interaction; confidentiality protects information exchanged during the
protocol. Each describes a distinct aspect of communication security.

\subsection{Policy-Based Authorization}

% Sources: paper/PAPER_STORY.md, Section 1;
% paper/sections/03_problem_formulation.tex, Policy--Authorization Consistency
% and Authorization Pipeline; README.md, Abstract and Agent Registration.
Policy-based access control expresses permissions through rules and applies
those rules when handling requests. Three concepts organize this process:
policy specification, policy evaluation, and enforcement.

\emph{Policy specification} defines the meaning of the rules, including which
subjects and interactions they cover and how applicable rules determine a
decision. In SAGA, a receiving user's contact rulebook associates agent-identity
patterns with budget values. Pattern matching identifies applicable rules;
the policy's rule-selection convention and interpretation of the selected
budget determine contact permission.

\emph{Policy evaluation} applies those semantics to a particular request.
Its inputs include the requesting identity, the target, and the applicable
policy. The result expresses whether the request is permitted under those
inputs. Policy meaning defines the decision to be computed, while the
evaluation procedure implements that computation.

\emph{Enforcement} applies the decision at the relevant protocol boundary.
In SAGA's communication flow, the Provider evaluates the receiving user's
stored rulebook before releasing the material needed to continue contact.
The receiving agent performs policy and credential checks, issues an
access-control token, and subsequently checks a presented token when handling
a message. These stages distinguish permission to proceed, credential
issuance, and message acceptance.

The authorization outcome thus depends on the specified policy semantics,
their evaluation for the request, and the application of the resulting
decision at the appropriate stages. A policy decision concerns a particular
initiating agent and receiving agent; where a credential mediates contact,
its use is part of this authorization context. This terminology provides the
basis for the policy--authorization consistency requirement developed next.
